\section{Introduction}

DBpedia extracts structured information from Wikipedia---infobox fields,
abstracts, categories, redirects and links between articles---and publishes it
as RDF, so that the content of an encyclopedia can be queried with SPARQL and
joined with other datasets~\cite{auer2007dbpedia, lehmann2015dbpedia}. The
English edition sits at the centre of the Linked Open Data cloud, and language
chapters publish the same kind of data for other Wikipedia
editions~\cite{kontokostas2012greek}. Vietnamese Wikipedia articles carry the
same semi-structured information: the infobox of a footballer lists his date of
birth, birthplace and every club he played for, season by season, but as
wikitext it cannot be queried. A question such as \emph{``Which players born in
Nghệ An have played for the national team, and for which club do they play
now?''} requires reading and joining several articles by hand.

This project addresses Topic~2 of the course, \emph{Build a DBpedia version for
Vietnamese language}, which consists of five tasks: (1)~define an ontology;
(2)~collect Vietnamese articles from Wikipedia; (3)~transform the data into
4-star Linked Data; (4)~establish links to the English DBpedia; and
(5)~provide an interface to query the data through a SPARQL endpoint or
terminal.

Rather than extracting all of Vietnamese Wikipedia shallowly, we extracted a
closed domain in depth: Vietnamese \emph{football} (players, clubs, national
teams, stadiums), \emph{provinces} and \emph{universities}. Its entities refer
to one another, so the graph supports multi-hop questions: player
$\rightarrow$ career station $\rightarrow$ club $\rightarrow$ home stadium
$\rightarrow$ province, player $\rightarrow$ birth province, university
$\rightarrow$ province. 777 of the 990 selected entities (78\%) have an English
Wikipedia article to link to, and 956 of the 990 articles (97\%) contain an
infobox, the main source that DBpedia extracts. Entities are selected through
Wikidata rather than through the category tree of Vietnamese Wikipedia, whose
categories editors assign inconsistently.

The contributions of this work are:
\begin{itemize}
  \item an OWL ontology, \term{vio:}, with 18 classes and 54 properties, aligned
        to the DBpedia Ontology for OWL~2~RL reasoning (Section~\ref{sec:ontology});
  \item a reproducible pipeline whose raw downloads are committed, so the
        dataset rebuilds offline (Sections~\ref{sec:collection}--\ref{sec:postprocess});
  \item 131{,}543 triples (88{,}898 asserted, 41{,}730 inferred) that pass eight
        quality checks with no error or warning (Section~\ref{sec:evaluation});
  \item 777 \term{owl:sameAs} links to English DBpedia, verified on the live
        endpoint, and 990 to Wikidata (Sections~\ref{sec:interlinking} and~\ref{sec:links});
  \item a server with a SPARQL~1.1 Protocol endpoint, IRIs dereferenced with 303
        and content negotiation, a terminal client, and a web interface with
        entity pages, an ontology browser, a SPARQL editor and step-by-step
        question answering (Section~\ref{sec:publishing}).
\end{itemize}

Table~\ref{tab:requirements} maps each task to its implementation; the following
sections cover background, design, publishing, evaluation and discussion.

\begin{table}[t]
\centering
\caption{Where each task of the assignment is implemented. Paths are under
\term{vidbpedia/}, except \term{ui/}.}
\label{tab:requirements}
\small
\begin{tabularx}{\textwidth}{@{}l>{\raggedright\arraybackslash}X>{\raggedright\arraybackslash}Xl@{}}
\toprule
Task & Implementation & Output & Section \\
\midrule
1. Ontology & \term{ontology/*.ttl}, merged by \term{kg/ontology.py}
  & 18 classes, 54 properties, 859 triples & \ref{sec:ontology} \\
2. Collect articles & \term{crawl/wikidata\_seeds.py}, \term{crawl/wiki\_enrich.py}
  & 990 entities with their Vietnamese articles & \ref{sec:collection} \\
3. 4-star data & \term{crawl/build\_rdf.py}, \term{kg/postprocess.py}
  & 131{,}543 triples in Turtle and N-Triples & \ref{sec:transform}, \ref{sec:postprocess} \\
4. Link to DBpedia & enwiki sitelinks of the Wikidata item
  & 777 \term{owl:sameAs dbr:} links, VoID linksets & \ref{sec:interlinking} \\
5. Query interface & \term{web/}: endpoint, Linked Data, Gradio; \term{ui/}: React; \term{kg/query.py}
  & \term{/sparql} endpoint, \term{query} command, Linked Data routes, web interface
  & \ref{sec:endpoint}, \ref{sec:deref}, \ref{sec:ui} \\
\bottomrule
\end{tabularx}
\end{table}
